% methods.tex (v10). Unnumbered subsections (Sci Rep). Supplementary Table S2 lists the years used in every analysis.
\section{Methods}\label{sec:methods}

\bmhead{Study regions and ALT labels}
ALT observations were taken from three public sources. For Alaska and north-western Canada we used the ABoVE soil moisture and active-layer thickness compilation, version 2\cite{Moore2025}, restricted to probe measurements made in August--October with $0<\mathrm{ALT}<300$ cm (measurements in January--July were excluded as early-season thaw depths), and the CALM end-of-season grid values\cite{Streletskiy2025}. For the Lena Delta we used the ALLena thaw-depth data\cite{Veremeeva2025}, restricted to August--September measurements without the 'exceeds probe length' flag, taking the seasonal maximum per location and year. Observations were assigned to 1 km cells and averaged within each cell and year, and the mean day of year of the measurements was kept; CALM values without a date were assigned day 244 (1 September). Only years with yearly climate input (2010--2024) were used. The resulting label set has \nncyAk{} cell-years in Alaska (\nyrsAk), \nncyLe{} in the Lena Delta (\nyrsLe) and \nncyCa{} in Canada (\nyrsCa), in \nnblkAk, \nnblkLe{} and \nnblkCa{} blocks of 0.5$^\circ$ (Table~1). \nnmultiAk, \nnmultiLe{} and \nnmultiCa{} cells have observations in four or more years.

\bmhead{Year-specific climate input}
Monthly ERA5-Land data for January 2010 to December 2024\cite{munozsabater2021_era5_land} (2 m air temperature, soil temperature of the top layer and snow water equivalent) were taken at the nearest land grid cell of each 1 km cell (fallback up to 0.5$^\circ$ for coastal cells). For each cell and year $y$ we computed eight climate variables: mean annual air temperature, TDD (sum of positive monthly mean temperatures times days), $\sqrt{\mathrm{TDD}}$, the warmest monthly mean, the freezing degree-days and the coldest monthly mean of the preceding winter (October $y-1$ to April $y$), the mean top-layer soil temperature and the mean snow water equivalent of the hydrological year (October $y-1$ to September $y$). For 2010 the winter variables use January--April only. The main input (year-specific climate input) accumulates TDD to the observation day by linear interpolation of the monthly cumulative TDD of the observation year. Two alternatives were run on the same cell-years: the full-season TDD of the observation year, and the 2011--2024 mean of the yearly values (mean climate input), which corresponds to the multi-year inputs of existing products. Maps use the full-season TDD of each year (end-of-season ALT).

\bmhead{Static covariates}
Terrain variables (elevation, slope, the sine and cosine of aspect, topographic position index and roughness) were computed from the 30 m Copernicus DEM\cite{CopDEM2021} in a 33 $\times$ 33 pixel window (about 1 km) around each cell. Soil variables (clay, sand, silt, bulk density, coarse fragments and pH at 5--15 cm, and soil organic carbon at 0--5, 5--15 and 15--30 cm) were sampled from SoilGrids 2.0 at 250 m\cite{Poggio2021}. The main feature set (23 variables) contains the eight climate variables, six terrain variables and nine soil variables. Sensitivity sets were climate only; climate and terrain; the main set plus MODIS summer NDVI and EVI and snow-off timing (MOD13Q1 and MOD10A2 version 6.1 for 2010--2024 via AppEEARS\cite{Didan2021,HallRiggs2021}); the main set plus high-resolution soil carbon and land-cover classes\cite{Palmtag2022}; and all variables including the climate of the previous year. The ESA CCI permafrost probability (1997--2021 mean)\cite{Westermann2024} was used only to define the map domain (probability $\geq$ 50\%), not as a model input.

\bmhead{Published products}
Six products were extracted at the label cells: Wei et al.\cite{Wei2026} (1 km, 2000--2024), Liu et al.\cite{Liu2024} (1 km, 2003--2020), ESA CCI Permafrost ALT v4 (1997--2021) and v5 (1997--2023)\cite{Westermann2024,Westermann2025}, Yi and Kimball\cite{YiKimball2019} (Alaska, 2001--2015) and Aalto et al.\cite{Aalto2018} (baseline 2000--2014, one value per cell). For yearly products the value of the observation year was used; the years available at the label cells were \nyrsprodWei{} (Wei et al.), \nyrsprodLiu{} (Liu et al.), \nyrsprodCcifour{} (CCI v4), \nyrsprodCcifive{} (CCI v5) and \nyrsprodYk{} (Yi and Kimball). Products were used only as comparison targets. Wei et al. publish the coordinates of their training sites; the distance from each label cell to the nearest training site was used for the 25 km mask.

\bmhead{Models}
The Stefan model writes $\mathrm{ALT}=E\sqrt{\mathrm{TDD}}$\cite{Stefan1891}. The source-coefficient Stefan model takes $E$ fitted by least squares on the other two regions; the recalibrated Stefan model fits $E$ on the training labels of the region. Direct ML predicts ALT from the 23 features with CatBoost gradient boosting\cite{Prokhorenkova2018} (600 iterations, learning rate 0.03, depth 6, L2 regularisation 3). Physics-input ML adds the recalibrated Stefan and Kudryavtsev\cite{Kudryavtsev1974} predictions as features. Physics pseudo-labels add to the training set an equal number of rows drawn from the cells of the scoring fold, with their covariates and the training-fold Stefan prediction as the target (no observed values). The physics-residual model keeps the recalibrated Stefan prediction $\hat{y}_P$ as the anchor and adds a CatBoost model $g$ of the residual:
\begin{equation}
\hat{y} = E\sqrt{\mathrm{TDD}} + \lambda\, g(\mathbf{x}), \qquad g \approx y - \hat{y}_P,
\end{equation}
with $\lambda = 0.5$ (0.25, 1 and a value chosen by inner block cross-validation were run as sensitivity cases). All learned models were run with three seeds and the predictions averaged. Random forest, a multilayer perceptron, TabM and FT-Transformer\cite{Gorishniy2021,Gorishniy2025,Holzmuller2024} replaced CatBoost in a learner sensitivity test. SHAP values of the residual and direct models were computed with the CatBoost implementation\cite{Lundberg2017} on models fitted to all labels of each region.

\bmhead{Validation designs}
Five designs were applied within each region\cite{Roberts2017,Ploton2020}: random five-fold by cell-year; spatial-block five-fold, with 0.5$^\circ$ blocks assigned to folds balanced by the number of cell-years (main design); year hold-out, in which each observation year with at least 20 cell-years is scored by models trained on the other years; block-by-year, in which both the block fold and the year of the scored cell-years are excluded from training; and region hold-out, in which a region is scored by models trained on the other two regions after removing training cells within 100 km of the test region. In region hold-out no target labels are available, so the recalibrated Stefan model equals the source-coefficient model.

\bmhead{Label ladders, transfer and external regions}
In the within-region ladder, the training labels of each spatial-block fold were subsampled to $n$ = 25, 50, 100, 250, 500, 1000, 2500 or all (three random draws), the Stefan coefficient was recalibrated on the same $n$ labels, and the full test fold was scored (all model structures in Supplementary Fig.~S11). In the transfer ladder, each region was scored in its spatial-block folds by models trained on all labels of the other two regions plus $n$ = 0, 25, 50, 100, 250 or 500 labels drawn from the training folds of the target region. The Stefan coefficient was refitted on the $n$ target labels ($n \geq 3$), and the residual model was trained on the source rows (with the source-coefficient anchor) and the target rows (with the refitted anchor). For the external test, CALM observations from 2010--2024\cite{Streletskiy2025} outside the three study regions and outside mountain permafrost (Switzerland, Kazakhstan) were assigned to 1 km cells. Censored values, measurements flagged as early, burnt, flooded or not reaching the permafrost table, and sites within 100 km of the study labels (Samoylov and Bykovsky in the Lena Delta) were removed, leaving \nextcy{} cell-years in \nextcells{} cells in eight regions (West, Central and North-East Siberia, European Russia, Chukotka, Kamchatka, Svalbard and Greenland). Climate inputs and covariates were computed in the same way. Each external region was scored block by block (cell by cell for Greenland, which has one block), with models trained on the three study regions plus $n$ = 0, 5, 10, 25, 50 or all labels from the other blocks of the same external region. One Svalbard borehole with a mean ALT of about 6 m, outside the training range, was excluded.

\bmhead{Scoring and statistics}
Root-mean-square error (RMSE), bias and Pearson correlation were computed on out-of-fold predictions. Differences between two methods ($\Delta$RMSE = RMSE of the first minus RMSE of the second) were computed on the same cell-years, and 95\% confidence intervals were taken from 2,000 bootstrap resamples of the 0.5$^\circ$ blocks (1,000 for the ladders). We describe a difference as resolved when its interval excludes zero. Conditions of the gain were examined within regional terciles of eight covariates (full-season TDD of the observation year, freezing degree-days, snow water equivalent, soil organic carbon, clay, slope, elevation and summer NDVI); terciles with fewer than 30 cell-years or fewer than five blocks were not scored. Year-wise RMSE was computed for years with at least 30 cell-years, and the correlation between observed and predicted year-to-year changes was computed at cells observed in several years.

\bmhead{Comparison with products}
Products and models were compared on the same cell-years. The common-cell comparison uses only cell-years at which Wei et al., Liu et al. and ESA CCI v4 all provide a value of the observation year (\nyrscc). A leakage mask removes cells within 25 km of the training sites of Wei et al.; Liu et al. were trained on an earlier version of the ABoVE data used here\cite{Liu2024}, so their Alaska and Canada values are not independent of our labels. In the product-anchored test, the product value of the observation year was linearly recalibrated on the training fold and corrected by the same residual learner ($\lambda$ = 0.5) on the same block folds.

\bmhead{Yearly maps, area of applicability and prediction intervals}
For Alaska and the Lena Delta the recalibrated Stefan and physics-residual models were fitted on all cell-years of the region and applied to the 1 km cells with an ESA CCI permafrost probability of at least 50\% (\nncellmapAk{} and \nncellmapLe{} cells), with the full-season climate of each year from 2011 to 2024. The area of applicability was defined following Meyer and Pebesma\cite{MeyerPebesma2021}: in the standardised space of the 23 features (with the 2011--2024 mean climate), a map cell is inside when its distance to the nearest training row does not exceed the 95th percentile of the distances from each training row to the nearest training row in another block. The 90\% prediction interval was taken from the 5th and 95th percentiles of the spatial-block out-of-fold residuals. For display, the 1 km values were interpolated bilinearly to the 250 m hillshade grid.

\bmhead{Computation}
Model fitting was run on cloud CPU nodes with 96 cores (Rescale) with fixed seeds, and data preparation, aggregation and figures on a local server. Supplementary Table S6 lists the runs.

\bmhead{Use of large language models}
A large language model (Claude, Anthropic) was used to assist with code development, figure scripting and language editing. All analyses, numbers and text were checked by the authors, and every number in the manuscript is generated from the stored result tables. [Authors to confirm.]

\bmhead{Code availability}
The code for data preparation, modelling, evaluation and figures is available at [repository and DOI]. Every number in the text is produced from the result tables by the scripts that generate the manuscript macros.
